Os experimentos foram planejados para responder às três hipóteses da
Seção~\ref{subsec:hipotese}. A Tabela~\ref{tab:plano_experimentos} relaciona cada
experimento à hipótese avaliada, aos dados utilizados e às métricas principais. Os
experimentos de reconhecimento foram executados no computador, sobre conjuntos de imagens
fixos e versionados no repositório. Os de latência, entrega e consumo foram executados com
o protótipo em bancada, com os firmwares dos dois fluxos descritos na
Seção~\ref{subsec:firmware}.

\begin{table}[H]
    \centering
    \caption{Plano de experimentação}
    \label{tab:plano_experimentos}
    \begin{tabularx}{\textwidth}{@{}>{\raggedright\arraybackslash}p{3.1cm}c>{\raggedright\arraybackslash}X>{\raggedright\arraybackslash}p{3.6cm}@{}}
        \toprule
        \textbf{Experimento} & \textbf{Hipótese} & \textbf{Dados e condição} & \textbf{Métricas} \\
        \midrule
        Efeito da quantização & H1 & Mesmo classificador em ponto flutuante e em INT8; 7.378
        recortes de validação e 1.894 imagens de teste independentes & Acurácia por recorte,
        leitura exata, acurácia por dígito, CER \\
        \addlinespace
        Leitura de ponta a ponta & --- & 39 fotos do conjunto de teste comum e 21 fotos de
        hidrômetros físicos; janela prevista pelo modelo e janela marcada à mão & Leitura exata,
        acurácia por dígito, CER, Dice e IoU \\
        \addlinespace
        Latência e entrega & H2 & Nó em ciclo autônomo com \textit{deep sleep}: 30 leituras no
        fluxo em nuvem e 34 no fluxo em borda; 30 envios adicionais de imagem com o nó sempre
        ligado & Latência ponta a ponta, \gls{pdr} por pacote, taxa de entrega por leitura \\
        \addlinespace
        Consumo de energia & H3 & Corrente do nó medida em 30 ciclos do fluxo em nuvem e 32
        do fluxo em borda & Carga por leitura (mAh), corrente acordado e em repouso, autonomia
        estimada \\
        \bottomrule
    \end{tabularx}
    \legend{Fonte: Elaborado pelos autores (2026).}
\end{table}

Como os dois fluxos usam os mesmos modelos INT8 e a mesma lógica de leitura
(Seção~\ref{subsec:visao_geral}), a acurácia do reconhecimento não é uma variável entre
eles. Por isso, a hipótese H1 é avaliada pela comparação entre as versões em ponto
flutuante e em INT8 do mesmo modelo, e não entre os fluxos. Já nos experimentos das
hipóteses H2 e H3, a única variável é o fluxo: o hardware, o enlace, o intervalo de repouso
e a bancada são os mesmos.

\subsection{Métricas}
\label{subsec:metricas_experimentos}

\subsubsection{Reconhecimento}
\label{subsubsec:metricas_ocr}

Para cada imagem, a leitura prevista $\hat{y}$ é comparada com a
leitura real $y$, ambas sequências de $N$ dígitos. São usadas três métricas sobre um
conjunto de $M$ imagens: a taxa de leitura exata, que conta apenas as leituras com todos os
dígitos corretos; a acurácia por dígito, que compara posição a posição; e a taxa de erro de
caracteres (\textit{character error rate}, CER), baseada na distância de edição de
Levenshtein $d_{\mathrm{L}}$, isto é, o número mínimo de inserções, remoções e substituições
que transformam uma sequência na outra:

\begin{equation}
    \label{eq:leitura_exata}
    Acc_{\mathrm{leitura}} = \frac{1}{M}\sum_{m=1}^{M} \mathbf{1}\!\left[\hat{y}_m = y_m\right]
    \qquad
    Acc_{\textrm{dígito}} = \frac{\sum_{m}\sum_{k=1}^{N_m} \mathbf{1}\!\left[\hat{y}_{m,k} = y_{m,k}\right]}{\sum_{m} N_m}
\end{equation}

\begin{equation}
    \label{eq:cer}
    CER = \frac{\sum_{m} d_{\mathrm{L}}(\hat{y}_m, y_m)}{\sum_{m} N_m}
\end{equation}

A taxa de leitura exata é a métrica mais exigente e a que interessa à aplicação: uma
leitura com um único dígito errado é uma leitura errada. Para o classificador isolado, a
métrica é a acurácia por recorte, a fração de recortes de um único dígito classificados
corretamente, acompanhada da acurácia por classe e da matriz de confusão.

A qualidade da localização da janela é medida pela sobreposição entre a máscara prevista
$P$ e a máscara de referência $G$, pelos coeficientes Dice e IoU (\textit{intersection over
union}):

\begin{equation}
    \label{eq:dice_iou}
    Dice = \frac{2\,|P \cap G|}{|P| + |G|}
    \qquad
    IoU = \frac{|P \cap G|}{|P \cup G|}
\end{equation}

\subsubsection{Latência e entrega}
\label{subsubsec:metricas_latencia}

A latência ponta a ponta de uma leitura é o intervalo entre
o disparo no nó e o instante em que o dado fica disponível no servidor: a imagem remontada,
no fluxo em nuvem, ou a leitura decodificada, no fluxo em borda. A entrega é avaliada de
duas formas. A primeira é o \gls{pdr} por pacote, definido na Equação~\eqref{eq:pdr}. A
segunda é a taxa de entrega por leitura, a fração das leituras disparadas cujo dado chegou
completo ao servidor:

\begin{equation}
    \label{eq:entrega_leitura}
    E_{\mathrm{leitura}} = \frac{\text{leituras disponíveis no servidor}}{\text{leituras disparadas}}
\end{equation}

As duas métricas diferem no fluxo em nuvem: como nenhum dos firmwares retransmite blocos no
modo autônomo, uma imagem com um único bloco perdido não é remontada, e a leitura se perde
mesmo com um \gls{pdr} alto. No fluxo em borda, em que cada leitura ocupa um pacote, as duas
coincidem. As taxas de entrega por leitura dos dois fluxos são comparadas pelo teste exato
de Fisher, adequado a proporções com amostras pequenas.

\subsubsection{Energia}
\label{subsubsec:metricas_energia}

A carga consumida em uma leitura é a integral da corrente $i(t)$ ao
longo de um ciclo completo, de um despertar ao seguinte, incluindo o repouso, calculada pela
regra do trapézio sobre as amostras e expressa em mAh:

\begin{equation}
    \label{eq:carga_leitura}
    Q_{\mathrm{leitura}} = \frac{1}{3600}\int_{t_{\mathrm{despertar},k}}^{t_{\mathrm{despertar},k+1}} i(t)\,dt
\end{equation}

\noindent com $i$ em mA e $t$ em segundos. Para estimar a autonomia em outros intervalos
entre leituras, o ciclo é modelado por dois estados: acordado, com corrente média
$I_{\mathrm{ac}}$ durante $t_{\mathrm{ac}}$, e em repouso, com corrente $I_{\mathrm{rep}}$
no restante do intervalo $T$. A autonomia $A$, em dias, para uma bateria de capacidade
$C_{\mathrm{bat}}$ em mAh, é então:

\begin{equation}
    \label{eq:autonomia_modelo}
    Q_{\mathrm{leitura}}(T) = \frac{I_{\mathrm{ac}}\,t_{\mathrm{ac}} + I_{\mathrm{rep}}\,(T - t_{\mathrm{ac}})}{3600}
    \qquad
    A = \frac{C_{\mathrm{bat}}}{Q_{\mathrm{leitura}}(T)} \cdot \frac{T}{86400}
\end{equation}

O modelo é ideal: não considera as perdas do regulador, a autodescarga nem a variação da
capacidade da bateria com a temperatura, de modo que a autonomia calculada é um limite
superior.

\subsection{Cenário de bancada}
\label{subsec:bancada}

Todas as medições em hardware foram feitas em 27 de outubro de 2026, em bancada, com o nó
sensor e o receptor a cerca de 2~m de distância. Nessa condição, o RSSI médio dos quadros
recebidos ficou entre $-34$ e $-27$~dBm, um sinal muito forte. Como antecipado na
Seção~\ref{subsec:simplificacoes}, o cenário avalia o comportamento dos fluxos, e não o
alcance do enlace. O enlace usou a configuração da Tabela~\ref{tab:parametros_enlace} nos
dois fluxos, e a câmera ficou fixa, apontada para o mostrador de um hidrômetro.

Nos experimentos de latência, entrega e consumo, os dois firmwares operaram no modo
autônomo com repouso de 30~s (\texttt{AUTO:30}), descrito na
Seção~\ref{subsubsec:ciclo_operacao}, e sem a janela para receber o comando de parada após
o despertar. O intervalo de 30~s não representa uma instalação real, em que as leituras
seriam espaçadas de minutos ou horas; ele foi escolhido para obter pelo menos 30 ciclos de
cada fluxo em poucas horas de medição. O efeito de intervalos mais longos é obtido pelo
modelo da Equação~\eqref{eq:autonomia_modelo}.

O nó foi alimentado pelos 5~V da porta USB da placa de gravação ESP32-CAM-MB, que também
fornecia a serial usada para acompanhar o log do firmware. O computador lia ao mesmo tempo
a serial do nó e a do receptor, de modo que todos os instantes foram marcados por um único
relógio, o do computador. A incerteza desses instantes, de alguns milissegundos, é
desprezível diante das latências medidas.

\subsection{Avaliação do reconhecimento}
\label{subsec:exp_ocr}

\subsubsection{Efeito da quantização (H1)}
\label{subsubsec:exp_quantizacao}

O classificador de dígitos treinado (Seção~\ref{subsubsec:classificador}) foi salvo em
ponto flutuante de 32 bits e, após a quantização (Seção~\ref{subsubsec:quantizacao_ocr}),
em INT8. As duas versões foram avaliadas sobre os mesmos 7.378 recortes de validação, da
divisão por foto descrita na Seção~\ref{subsubsec:dados_ocr}, por um \textit{script} que
apenas carrega os modelos já treinados e mede, sem treinar novamente. Como os pesos de
origem e os dados são os mesmos, a diferença de acurácia entre as versões isola o efeito
da quantização.

Como os recortes de validação vêm das mesmas fontes usadas no treino, o modelo INT8 também
foi avaliado nas pastas de teste do \textit{Word-Wheel Water Meter Dataset}, que ficaram
fora de todo o treinamento: 1.894 imagens da janela de dígitos, de $200 \times 60$ pixels.
Cada imagem foi dividida em $N$ fatias iguais, como no firmware, e cada fatia foi
classificada. Como essas imagens já são recortes retos da janela, esse teste mede o
classificador e o fatiamento, sem a etapa de localização. O teste foi repetido com o modelo
ao fim da fase A e com o modelo final, após o ajuste fino da fase B, para verificar se o
ajuste às imagens do \gls{esp32cam} prejudicou a generalização.

\subsubsection{Leitura de ponta a ponta}
\label{subsubsec:exp_leitura}

Para avaliar o pipeline completo, da imagem da câmera até a leitura, foram usados dois
conjuntos de fotos capturadas pelo \gls{esp32cam}, nenhum deles usado no treino ou no
ajuste dos modelos.

O \textbf{conjunto de teste comum} reúne 39 fotos de um monitor que exibia imagens de
hidrômetros, com 19 leituras distintas de dois modelos de medidor, um de cinco e outro de
seis roletes. As fotos foram selecionadas entre 152 capturas, com dois critérios: sem
rotação de 90° ou 180° e sem nenhum rolete parado entre dois números, caso em que a
leitura real é ambígua. Por comparação de \textit{hash}, verificou-se que nenhuma delas é
idêntica às fotos usadas no treino. O gabarito de cada foto foi conferido manualmente,
sobre uma folha que imprime a leitura esperada sobre a imagem, e para cada foto também foi
marcada à mão a janela de dígitos.

O \textbf{conjunto de hidrômetros físicos} reúne 21 fotos de medidores reais, de
diferentes modelos, com cinco a oito roletes, capturadas pelo \gls{esp32cam} no quadro de
$160 \times 120$ pixels. Ao contrário do conjunto comum, essas fotos têm reflexos no vidro,
iluminação variada e ângulos diversos. A leitura de cada foto foi anotada manualmente. Para
cada foto, foram marcados também a inclinação da fileira de roletes, a janela de dígitos no
quadro endireitado e a polaridade da janela: dígitos escuros em fundo claro ou dígitos
claros em fundo escuro. A inclinação foi a estimada pela transformada de Hough quando ela
deixava a fileira reta, em 16 fotos, e corrigida à mão nas outras cinco. A janela foi marcada
pela geometria da imagem, sobre o quadro ampliado com uma grade de coordenadas, a partir das
bordas da janela e dos centros dos roletes, de modo que as divisões das fatias caíssem entre
eles. Todas as marcações foram concluídas antes da execução da leitura e não foram ajustadas
pelo resultado.

Cada foto foi avaliada nos dois fluxos. No fluxo em borda, o quadro de
$160 \times 120$ pixels é lido diretamente; no fluxo em nuvem, o mesmo quadro é antes
codificado em JPEG com qualidade 90 e decodificado, como ocorre na transmissão. Essa
avaliação é uma emulação em Python, com os mesmos arquivos \texttt{.tflite} INT8, e não a
execução na placa. No conjunto comum, as fotos foram lidas em três configurações:

\begin{enumerate}
    \item \textbf{Janela prevista pelo modelo}: o pipeline completo, como no firmware, com
    a caixa obtida da máscara do modelo de localização.
    \item \textbf{Janela marcada à mão}: a caixa é informada diretamente e o modelo de
    localização não é usado. Essa configuração corresponde à ROI calibrada na instalação
    (Seção~\ref{subsubsec:leitura_roi}) e mede o classificador e o fatiamento quando a
    janela está correta.
    \item \textbf{Janela marcada à mão com correção de inclinação}: o quadro é antes
    endireitado no servidor pelo ângulo estimado com a transformada de Hough, recurso
    disponível apenas no fluxo em nuvem (Seção~\ref{subsubsec:recursos_servidor}).
\end{enumerate}

A distinção entre essas configurações é essencial para interpretar os resultados.
Quando a janela é prevista pelo modelo, o erro da leitura é dominado pelo erro da
localização: com uma caixa que cobre quase todo o mostrador, as fatias não correspondem
aos roletes, e nenhum classificador acertaria a leitura. Assim, os resultados da primeira
configuração medem o modelo de localização, e não a capacidade de leitura do classificador,
que é medida pelas configurações com a janela marcada à mão.

O conjunto de hidrômetros físicos, cujas fotos têm inclinações de até cerca de 55°, foi lido
com a janela prevista pelo modelo e com a janela e a inclinação marcadas à mão. Nesse caso,
a rotação do quadro foi aplicada nos dois fluxos, o que no firmware exigiria uma rotação com
ângulo fixo, ainda não implementada. Como variante de diagnóstico, as fotos de fundo escuro
foram lidas também com os tons da janela invertidos, o que equivale a informar a polaridade
do medidor na instalação, assim como o número de roletes. Depois da avaliação do conjunto
completo, os resultados também foram calculados para o subconjunto de fotos que atendem às
condições de uma instalação adequada: câmera frontal, dígitos escuros em fundo claro e todos
os dígitos visíveis no quadro. Por ter sido definido após a avaliação, esse subconjunto é
apresentado como complemento, e o conjunto completo continua sendo a referência.

Por fim, como a ROI de uma instalação real é calibrada uma única vez e pode ter pequenos
desvios, foi medida a tolerância da leitura ao erro de posição da janela. Partindo da janela
marcada à mão nas 39 fotos do conjunto comum, já endireitadas e no fluxo em nuvem, cada uma
das quatro coordenadas da caixa recebeu um erro inteiro sorteado uniformemente entre $-d$ e
$+d$ pixels, para $d \in \{0, 1, 2, 3, 4, 6\}$, com três sementes aleatórias. Cada caso foi
lido com fatias iguais e com a votação entre cinco caixas deslocadas.

\subsection{Latência e entrega (H2)}
\label{subsec:exp_latencia}

A latência e a entrega foram medidas com o \textit{script} \texttt{latencia\_ponta\_a\_ponta.py},
com o nó no modo autônomo. Nesse modo, o computador apenas escuta: o disparo é a primeira
linha que o firmware escreve na serial após o despertar, e o fim é o instante em que a
imagem é remontada ou a leitura é decodificada. Como o boot da ROM do ESP32 não aparece na
serial, o disparo marcado ocorre cerca de 0,3 a 0,5~s após o despertar real, o que
subestima a latência nesse valor nos dois fluxos. No fluxo em nuvem, a latência não inclui a
execução do \gls{ocr} no servidor, que é feita depois sobre as imagens salvas; como mostra a
Tabela~\ref{tab:tempo_inferencia}, esse tempo é da ordem de milissegundos e não altera a
comparação. Foram feitas duas corridas:

\begin{itemize}
    \item \textbf{Fluxo em nuvem}: 31 despertares, das 9h34 às 11h20, simultâneos à medição
    de corrente. Em 30 deles o nó enviou uma imagem; no outro, o nó acordou e voltou a dormir
    sem enviar nada, e esse despertar foi excluído.
    \item \textbf{Fluxo em borda}: 34 despertares, das 12h59 às 13h27.
\end{itemize}

Para conferir o \gls{pdr} do fluxo em nuvem com uma contagem exata dos pacotes enviados, foi
feita uma terceira corrida com o \textit{script} \texttt{reconstrutor.py} no modo controlado.
Nele, o nó fica sempre ligado, sem \textit{deep sleep}, e o computador envia o comando de
captura, espera o fim do envio e, após uma pausa de 10~s, dispara a leitura seguinte. Foram
30 envios, das 14h59 às 16h44. O número de pacotes enviados vem do log do próprio nó, que
registra cada tentativa de transmissão, e não da lacuna no contador de quadros (FCnt), que
não detecta perdas no início nem no fim de uma imagem. Cada corrida grava todos os quadros
recebidos, os logs das duas seriais e as imagens remontadas, que ficam versionadas no
repositório. O tempo no ar de cada envio é calculado pela Equação~\eqref{eq:tpacote_enlace},
e não medido.

Em algumas execuções, o receptor deixou de repassar quadros ao computador por 15 a 25
minutos seguidos, enquanto o nó registrava todos os envios como aceitos pelo módulo. Isso
ocorreu em uma primeira tentativa do fluxo em borda, das 12h25 às 12h46, e em cinco execuções
do fluxo em nuvem, das 13h41 às 14h55. Essas execuções foram interrompidas e mantidas no
repositório em arquivos separados, como registro, mas não entram nas estatísticas: as perdas
nelas vêm de uma falha do receptor, cuja causa não foi identificada, e não do enlace. Essa
exclusão é retomada na discussão das limitações.

\subsection{Consumo de energia (H3)}
\label{subsec:exp_energia}

A corrente do nó foi medida com um sensor INA219, cujo resistor \textit{shunt} de
0,1~$\Omega$ foi ligado em série na alimentação de 5~V, antes da \gls{esp32cam}. Como o
Radioenge é alimentado pelo regulador de 3,3~V da própria placa, a medição inclui o nó
inteiro: microcontrolador, câmera, módulo de rádio, regulador e LED. Um capacitor de
100~µF, ligado depois do resistor, evita a queda de tensão nos picos de transmissão. O
INA219 foi lido por I2C por um Arduino Uno, que enviava ao computador uma amostra a cada
10~ms (100~Hz), com a tensão no resistor e a tensão de alimentação. A
Figura~\ref{fig:bancada_corrente} mostra a montagem.

\begin{figure}[H]
    \centering
    \caption{Montagem da medição de corrente do nó sensor}
    \label{fig:bancada_corrente}
    \begin{tikzpicture}[
        font=\footnotesize,
        node distance=8mm and 10mm,
        bloco/.style={draw, rounded corners=2pt, align=center, fill=white,
                      minimum height=9mm, text width=#1},
        bloco/.default=3.2cm,
        seta/.style={-{Stealth[length=2mm]}, thick},
        rotulo/.style={font=\scriptsize\itshape, align=center}
    ]
        \node[bloco] (mb) {\textbf{ESP32-CAM-MB}\\5~V da USB e serial do nó};
        \node[bloco, right=of mb] (ina) {\textbf{INA219}\\\textit{shunt} de 0,1~$\Omega$};
        \node[bloco, right=of ina] (cam) {\textbf{\gls{esp32cam}}\\regulador de 3,3~V};
        \node[bloco, below=of cam] (radio) {\textbf{Radioenge RD49C}};
        \node[bloco, below=of ina] (uno) {\textbf{Arduino Uno}\\100 amostras/s};
        \node[bloco, below=of mb] (pc) {\textbf{Computador}\\\texttt{corrente.py}};
        \node[draw, dashed, rounded corners=4pt, inner sep=2.5mm, fit=(cam) (radio),
              label={[font=\scriptsize\bfseries, anchor=south east]north east:Nó sensor medido}] {};

        \draw[seta] (mb) -- node[rotulo, above] {5~V} (ina);
        \draw[seta] (ina) -- node[rotulo, above] {5~V} (cam);
        \draw[seta] (cam) -- node[rotulo, right] {3,3~V e UART} (radio);
        \draw[seta] (ina) -- node[rotulo, right] {I2C} (uno);
        \draw[seta] (uno) -- node[rotulo, above] {USB} (pc);
        \draw[seta] (mb) -- node[rotulo, left] {USB} (pc);
    \end{tikzpicture}
    \legend{Fonte: Elaborado pelos autores (2026).}
\end{figure}

O INA219 foi configurado com o ganho que admite até 320~mV no resistor (3,2~A), para não
saturar nos picos de transmissão. A corrente é calculada no computador pela tensão no
resistor, cujo passo de 10~µV corresponde a 0,1~mA. Com o nó desconectado, a leitura foi de
0,00~mV, de modo que nenhum \textit{offset} foi descontado. Considerando a tolerância do
resistor e o \textit{offset} típico do sensor, a incerteza é de cerca de 1~mA mais 1\% da
leitura, o que torna confiáveis apenas as correntes de repouso acima de alguns
miliampères. A taxa de 100 amostras por segundo registra cerca de 13 amostras em cada
pacote de 135,8~ms no ar.

A análise, feita pelo \textit{script} \texttt{corrente.py}, separa os ciclos pelos trechos
de repouso. O nível de repouso é estimado pela mediana dos 10\% menores valores da
captura, e todo trecho abaixo de um limiar, de 5~mA ou 50\% acima desse nível, o que for
maior, com pelo menos 12~s de duração, é considerado repouso. Cada ciclo vai do fim de um
repouso ao fim do seguinte, e a carga do ciclo é calculada pela
Equação~\eqref{eq:carga_leitura}. A duração mínima de 12~s foi necessária porque picos
isolados de uma única amostra, de cerca de 61~mA, partiam alguns trechos de repouso em
dois e criavam ciclos falsos. Foram feitas duas capturas principais:

\begin{itemize}
    \item \textbf{Fluxo em nuvem}: 685.000 amostras em 6.850~s, a partir das 9h29. Dos 31
    ciclos encontrados, um foi excluído: o mesmo despertar sem envio excluído na medição de
    latência, em que o nó ficou acordado por 0,97~s. Restaram 30 ciclos completos.
    \item \textbf{Fluxo em borda}: 169.738 amostras em 1.697~s, das 12h22 às 12h49, com 32
    ciclos completos.
\end{itemize}

Antes dessas capturas, um teste curto de quatro ciclos com uma versão anterior do firmware
do fluxo em nuvem, que mantinha a câmera ligada durante todo o envio, foi usado como
referência para medir o efeito de desligá-la após a captura.

A partir das médias medidas de corrente e tempo acordado de cada fluxo, a autonomia foi
estimada pela Equação~\eqref{eq:autonomia_modelo} para uma bateria hipotética de
2.000~mAh e para intervalos de 10~minutos e de 1~hora entre leituras. A corrente de
repouso medida, de cerca de 24~mA, inclui o Radioenge, que permanece ligado durante o
\textit{deep sleep} do ESP32 (Seção~\ref{subsubsec:ciclo_operacao}), e os demais
componentes da placa. A medição que separaria essas parcelas, com o módulo desconectado,
não foi realizada. Por isso, a estimativa foi repetida com correntes de repouso hipotéticas
de 5, 1 e 0,1~mA, para mostrar como o resultado depende desse valor.
